\documentclass[10pt,letterpaper]{article}
\usepackage{cornell-notes}

\title{Chapter 12: Random Set Partition}
\author{}
\date{}

\setDocTitle{Chapter 12: Random Set Partition}
\setDocSubtitle{Combinatorial Algorithms}
\setDocVersion{v1.0}
\setDocStatus{Study Companion}
\setDocOwner{Combinatorial Algorithms Cornell Notes}
\setDocAuthor{}
\setDocDate{\today}
\setDocSummary{Cornell-format study and review notes for Chapter 12: Random Set Partition.}

\setCornellCollection{Combinatorial Algorithms Cornell Notes}
\setCornellUnitType{Chapter}
\setCornellUnitNumber{12}
\setCornellUnitTitle{Random Set Partition}
\setCornellCompanionLabel{Study and review companion}

\hypersetup{pdftitle={Chapter 12: Random Set Partition},pdfsubject={Cornell notes},pdfauthor={}}

\begin{document}
\maketitle
\begin{tcolorbox}[colback=Paper,colframe=Ink]{\small\bfseries\color{Teal}CORNELL NOTES}\par{\LARGE\bfseries Chapter 12: Random Partition of an $n$-Set (RANEQU)}\par\medskip\textbf{Topic:} Bell-number recurrence sampling\hfill\textbf{Date:} \today\par\textbf{Study purpose:} Sample class sizes, then randomize labels.\end{tcolorbox}
\begin{tcolorbox}[colback=Teal!5,colframe=Teal,title=Essential question]How does the Bell-number recurrence determine class-size probabilities for a uniformly random set partition?\end{tcolorbox}
\begin{tcolorbox}[colback=white,colframe=Mist,title=Learning objectives]\begin{itemize}\item Derive the recurrence for Bell numbers.\item Interpret a recurrence term as the size of the class containing the largest current label.\item Build class labels without relabeling intermediate sets.\item Explain why a final random permutation restores uniform labeled membership.\item Identify the scaled count table used by the implementation.\end{itemize}\textbf{Prerequisites:} set partitions, binomial coefficients, induction, random permutations.\end{tcolorbox}
\section*{Cornell notes}
\begin{longtable}{@{}Q!{\color{Mist}\vrule width 1pt}A@{}}\cornellhead
\cue{What are the counts?}{Let $a_n$ be the number of partitions of an $n$-element set, the Bell number. By selecting the $k$ elements outside the class containing element $n$, then partitioning those $k$ elements,
\[
a_n=\sum_{k=0}^{n-1}\binom{n-1}{k}a_k,
\qquad a_0=1.
\]}
\cue{What probability follows?}{At residual size $m$, the distinguished class has $m-k$ elements with probability
\[
\Pr(k)=\frac{\binom{m-1}{k}a_k}{a_m},\qquad 0\le k<m.
\]
The numerator counts partitions whose class containing the distinguished last element leaves exactly $k$ elements for further partitioning.}
\cue{How are class sizes stored?}{Choose a class size, write one class number into a contiguous block at the end of the $q$ array, reduce $m$, increment the class number, and repeat. The sizes themselves need not be retained because their labels are written immediately.}
\cue{Why apply a random permutation?}{The block construction privileges positions. Randomly permuting the completed label array assigns the chosen class sizes uniformly to the actual element labels. The output $q_i$ is then the class containing element $i$.}
\cue{Why is it uniform?}{Fix a partition with classes $T_1,\ldots,T_\ell$. Sum over which class is selected first. The class-size probability, the probability that the shuffled block occupies exactly that class, and the inductive probability of the remaining partition multiply; factorial and Bell-number factors cancel. Summing over all first classes gives $1/a_n$.}
\cue{What is actually tabulated?}{To control numeric magnitude, the implementation stores $b_n=a_n/n!$ rather than $a_n$. A nested evaluation of the recurrence avoids repeatedly forming large binomial coefficients and factorials.}
\cue{Routine relationships}{RANEQU uses a random-number generator for weighted class-size choice and calls RANPER for the final shuffle. The output includes the number of classes and the class-label vector.}
\cue{Cautions}{Class labels are arbitrary after shuffling; compare partitions by equivalence classes, not by raw label numbers. Count scaling reduces but does not eliminate floating-point concerns. A distribution test should aggregate canonicalized partitions.}
\cue{Retrieval practice}{\begin{enumerate}\item What does $a_n$ count?\item Which class motivates the recurrence?\item Why are contiguous blocks sufficient?\item Why shuffle?\item What scaled quantity is stored?\end{enumerate}}
\end{longtable}
\begin{tcolorbox}[colback=Ink!4,colframe=Ink,title=Cornell summary]
RANEQU obtains a uniform set partition from the Bell-number recurrence. The recurrence classifies partitions by the number of elements left outside the class containing the current distinguished element. After normalization, its terms give valid probabilities for choosing the next class size. The routine writes class labels into blocks while residual elements remain, then randomly permutes the label array so the blocks are assigned uniformly to actual elements. An induction over the first selected class proves that every partition receives probability $1/a_n$. The implementation stores scaled Bell numbers $a_n/n!$ and delegates the last randomization to the permutation sampler.
\end{tcolorbox}
\end{document}
